\section{Quadratic orbits and exact obstructions}
\label{fd:geometry-section}

Quadratic sets can be brought to a difference of squared norms by
homogenization and a linear change of coordinates. This does not identify
a unique free set: different coordinates can give different cuts. We
first describe a signature for which the unrestricted transformation
orbit is complete, and then the bilinear signature for which it is not.

\subsection{One negative or one positive square}

Let $Q_{n,m}=\{(x,y)\in\R^n\times\R^m:\norm{x}\le\norm{y}\}$ and
$C_\ell=\{(x,y):\ell^Tx\ge\norm{y}\}$ for a unit vector $\ell\in\R^n$.
An automorphism preserves the homogeneous quadratic form
$\norm{x}^2-\norm{y}^2$. This geometry is related to the general maximal
quadratic-free characterizations of \citet{MPS2025,MPS2026}.

\begin{theorem}[Complete Lorentz orbits]\label{fd:lorentz}
For $n\ge2$, every full-dimensional maximal $Q_{n,1}$-free set is
\begin{equation}\label{fd:lorentz-pair}
 C_{\gamma_1,\gamma_2}
 =\{(x,y):\gamma_1^Tx\ge y,\ \gamma_2^Tx\ge-y\},
 \quad \norm{\gamma_i}=1,\quad\gamma_1\ne-\gamma_2.
\end{equation}
For any fixed unit $\ell$, these are exactly the images of $C_\ell$
under $SO^+(n,1)$, the determinant-one time-orientation-preserving
Lorentz group. For signature $(1,m)$, the only full-dimensional maximal
free sets are $\{x\ge\norm{y}\}$ and $\{-x\ge\norm{y}\}$.

If a quadratic's homogenized nondegenerate part has signature $(n,1)$
or $(1,m)$, every full-dimensional maximal free affine set is a slice
of one of these homogeneous sets, with the radical coordinates included
as a product before slicing. Thus the free-direction transformation orbit
has single-cut supremum $\zK(\rc)$ for positive costs.
\end{theorem}
\Cref{fd:lorentz-proof} gives elementary separation, classification,
Lorentz-basis, and slicing proofs. In particular no transitivity theorem
is needed. Every indefinite quadratic in two variables has one of these
homogenized signatures. Equality of the supremum does not imply
attainment: the positive-cost example in \cref{fd:boundary-examples}
already lies in this class.

For signature $(n,1)$, testing whether the target simplex fits inside
\cref{fd:lorentz-pair} is two independent second-order cone feasibility
problems in $\gamma_1,\gamma_2$. The relaxation
$\norm{\gamma_i}\le1$ remains free; the classification proof extends any
such full-dimensional pair to a maximal unit pair. Hence the relaxation
preserves the best value. Signature $(1,m)$ needs only two candidates.

The implemented point rule is a different restriction. In transformed
coordinates it chooses $\ell=x'/\norm{x'}$, where $(x',y')$ is the
homogenized apex, rather than optimizing the free direction.
\begin{proposition}[The transformed point rule]\label{fd:point-rule}
For signature $(n,1)$, at a fixed apex $u=(x,y)$ outside $Q_{n,1}$,
the plain point rule under all transformations produces exactly the
pairs in \cref{fd:lorentz-pair} for which
\begin{equation}\label{fd:point-span}
 u=a(\gamma_1,1)+b(\gamma_2,-1),\qquad a,b>0.
\end{equation}
For $n=2$ this is a one-parameter family. Writing
$\gamma_1=(\cos\theta,\sin\theta)$, its second direction is determined by
\[
 a=\frac{\norm{x}^2-y^2}{2(\gamma_1^Tx-y)},\qquad
 b=a-y,\qquad \gamma_2=\frac{x-a\gamma_1}{b},
\]
on the domain where $a,b>0$. The unrestricted pairs have two parameters.
\end{proposition}
The proof is in \cref{fd:point-rule-proof}. In general the respective
dimensions are $n-1$ and $2(n-1)$. When both null tangency rays meet an
affine slice with positive scaling, \cref{fd:point-span} says that the
apex lies on the chord between their two finite tangency points.
There are also maximal affine sets whose contacts are only asymptotic:
for $S=\{(x,y):xy\ge1\}$, the quadrant $\{x\ge0,y\le0\}$ is maximal
free and misses $S$ entirely. Its two homogeneous tangency rays lie at
infinity, so no affine apex satisfies \cref{fd:point-span}. Neither the
plain point rule nor an enlargement of one of its sets produces this
quadrant at any apex. \Cref{fd:asymptotic-example} supplies the exact proof.
This exclusion concerns generation of sets; it alone asserts no strict
gap for the point-rule supremum at a particular corner.

Even a fixed eigen-coordinate choice can lose almost the entire bound.
\begin{proposition}[A fixed rule with vanishing ratio]\label{fd:fixed-rule}
Let $S=\{(x,y):y^2\ge x^2+1\}$, $\sbar=(x_0,0)$, rays
$(-1,0),(0,1)$, and costs $(\epsilon,1)$, where
$0<\epsilon<1$ and $x_0\ge\epsilon/\sqrt{1-\epsilon^2}$.
The eigen-coordinate point rule gives
\[
 \left\{(x,y):|y|\le\frac{x_0x+1}{\sqrt{1+x_0^2}}\right\},\qquad
 z_{\rm rule}=\min\left\{\epsilon(x_0+1/x_0),\sqrt{1+x_0^2}\right\}.
\]
The strip $|y|\le1$ belongs to the same free-direction family and has
value one, whereas
$\zK=\epsilon x_0+\sqrt{1-\epsilon^2}$.
With $\epsilon=x_0^{-2}$ and $x_0\to\infty$,
$z_{\rm rule}/\zK\to0$ and the strip's ratio tends to one.
\end{proposition}
\Cref{fd:fixed-rule-proof} proves the formulas.
Even the closure of the entire constant-direction family can have a
vanishing ratio on pointed corners; \cref{fd:constant-family} proves
that separate assertion. Neither statement concerns the unrestricted
transformation orbit in \cref{fd:lorentz}.

\subsection{The determinant model for a bilinear constraint}

An indefinite quadratic of rank two is a product of two independent
linear forms plus an affine term. If its linear term belongs to the
range of its quadratic matrix, only two affine coordinates matter and
\cref{fd:lorentz} applies. Otherwise a third independent affine coordinate
puts it in the form $w-xy$ or its negative. We use
\begin{equation}\label{fd:det-model}
 q(s)=w-xy,\quad S=\{q\le0\},\quad
 M(s)=\begin{pmatrix}w&x\\y&1\end{pmatrix},\quad
 M_0(p)=\begin{pmatrix}p_w&p_x\\p_y&0\end{pmatrix},\quad
 E=\begin{pmatrix}1&0\\0&0\end{pmatrix}.
\end{equation}
The homogeneous form $\det M$ has signature $(2,2)$, and $E$ increases
the $w$ coordinate. The opposite bilinear side follows by
$(x,y,w)\mapsto(-x,y,-w)$.
Throughout the bilinear corner statements, $q(\sbar)>0$ and $\rc>0$.
For this $S$, the cone-containment hypothesis in
\cref{fd:cone-attainment} holds exactly when $\cone(P)=\R^3$.
Indeed every $(x,y,w)$ is the midpoint of
$(x+a,y+a,w)$ and $(x-a,y-a,w)$, both in $S$ for sufficiently large
$a$, so $\conv S=\R^3$. Any convex cone containing $S-\sbar$ must
therefore be the whole space. The opposite bilinear side follows
by the stated symmetry.

\begin{proposition}[The bilinear orbit]\label{fd:orbit}
The determinant-preserving linear maps on $2\times2$ matrices are
$M\mapsto AMB^T$ and $M\mapsto AM^TB^T$, with
$\det A\det B=1$. The homogeneous constant-direction free cone is,
in suitable coordinates, $\{M:\sym(M)\succeq0\}$. Its full orbit is
\[
 \widehat C_F=\{M:\sym(F^TM)\succeq0\},\qquad\det F>0,
\]
modulo positive scaling of $F$. Its affine slice is $C_F$.
These cones are full dimensional and determinant-free. At a nonzero
rank-one matrix $M=ab^T$, membership is exactly
$F^Ta\in\R_+b$. On the slice, their feasible contacts form the increasing
Möbius graph
\[
 y=\frac{F_{11}x+F_{21}}{F_{12}x+F_{22}},\qquad
 F_{12}x+F_{22}>0,\qquad w=xy.
\]
\end{proposition}
The matrix identities and the automorphism classification are proved in
\cref{fd:orbit-proof}. These are three-parameter sets. The one-parameter
family (14a) of \citet{BCM2020} is exactly the rotation subfamily
$F^T\in SO(2)$; \cref{fd:bcm} gives the identities and shows why
nonrotation orbit members lie outside both of their displayed families.

\begin{proposition}[Upward completion]\label{fd:completion}
For $\det F>0$, define
\[
 B_F=\cl(C_F+\R_+e_w),\qquad Z_F=\sym(F^TE).
\]
The corresponding homogeneous completion retains exactly the inequalities
$\ip{Fvv^T}{M}\ge0$ for which $v^TZ_Fv\ge0$. On the affine slice,
\begin{equation}\label{fd:lowering}
 s\in B_F\quad\Longleftrightarrow\quad
 \exists\tau\in[0,q(s)]:\sym(F^T[M(s)-\tau E])\succeq0.
\end{equation}
Writing $F=\bigl(\begin{smallmatrix}a&b\\c&d\end{smallmatrix}\bigr)$,
the exceptional class $b=0,a<0$ has empty completed slice. Every other
$B_F$ is full dimensional, maximal $S$-free, and contains $C_F$.
\end{proposition}
The proof in \cref{fd:completion-proof} uses an exact two-term PSD
decomposition and treats closure on the slice explicitly. These
completions agree with the applicable maximal completion construction
of \citet{MPS2026}. The constant-direction member in the eigen-coordinates
of \citet{ChmielaMunozSerrano2020ZIB} is their bilinear Case~4 set: its defining function
is the support function of the retained spherical cap.

Let $\famA$ consist of the orbit slices $C_F$ that contain $\sbar$ in
their interior, and let $\famB$ consist of the nonempty completions
$B_F$ with that property. Write $\zA,\zB$ for the corresponding
single-cut suprema. Then $\zA\le\zB\le\zK$. A completion can contain
the apex even when its underlying orbit slice does not.

\begin{proposition}[Convex feasibility for the best orbit value]
\label{fd:orbit-sdp}
For $t>0$, consider the nested feasibility system
\[
 \sym(F^TM(\sbar))\succ0,\qquad
 \sym\left(F^TM\left(\sbar+\frac{t}{\omega_j}p_j\right)\right)
 \succeq0\quad(1\le j\le N).
\]
Feasibility at target $t$ is equivalent to the existence of an orbit cut
of value at least $t$. The strict apex condition implies $\det F>0$ and,
by positive scaling, may be replaced by
$\sym(F^TM(\sbar))\succeq I$. Thus bisection uses $N+1$ linear matrix inequalities (LMIs) of size
two in four scalar variables.
\end{proposition}
\begin{proof}
The endpoint LMIs express $T_t\subseteq C_F$, and the apex LMI expresses
interior membership. A positive definite symmetric part has positive
determinant, so $\det(F^TM(\sbar))>0$; since $q(\sbar)>0$, this implies
$\det F>0$. Scaling normalizes the smallest apex eigenvalue. Convexity
and nestedness follow directly from the affine LMIs and the apex-to-endpoint
segments. The supremum need not itself be feasible.
\end{proof}

\subsection{A tangent edge forces a pencil}

Let $t$ be a feasible boundary point in the relative interior of a
nontrivial segment $e$, with direction $d$ satisfying
$\nabla q(t)^Td=0$ and $\det M_0(d)>0$. Such a segment occurs when a
two-ray minimizer lies in the relative interior of its objective face.
Orient $d$ so that $d_x>0$, and set
$J=\bigl(\begin{smallmatrix}0&1\\-1&0\end{smallmatrix}\bigr)$.

\begin{lemma}[Tangent-edge rigidity]\label{fd:tangent-pencil}
If an invertible $F$ has $e\subseteq B_F$, then
\begin{equation}\label{fd:pencil}
 F^T=\theta\bigl(G_0+uG_1\bigr),\qquad
 G_0=J\adj(M_0(d)),\quad G_1=J\adj(M(t)),\quad
 \theta>0,\quad u\in\R.
\end{equation}
The pencil has constant determinant $\det M_0(d)>0$ and
$\sym(G_1M(t))=0$. If the nontrivial segment lies on a ruling, so that
$\det M_0(d)=0$, no invertible $F$ has $e\subseteq B_F$.
\end{lemma}
\Cref{fd:pencil-proof} derives the pencil from the two-sided PSD
conditions. Without the chosen orientation the sign condition would
be $\theta d_x>0$. A tangent edge therefore reduces the orbit-containment
test to one real parameter, but the remaining simplex vertices can impose
incompatible conditions on it.

\begin{theorem}[A strict gap for both transformation families]
\label{fd:counterexample}
For $S=\{w\le xy\}$, take $\rc=(1,1,1)$ and rays from
\[
 \sbar=(-9/2,0,3/2)
\]
to
\[
 v_1=(-1,-6,18),\qquad v_2=(-5,6,-18),\qquad v_3=(0,5/2,5/2).
\]
Then $\zK=1$, attained only at
$t=(-3,0,0)=(v_1+v_2)/2$. Neither an orbit slice nor an upward completion
contains $T_1$, and
\[
 \zA\le\zB<\zK.
\]
The unrestricted optimum is attained by a maximal free set outside
$\famB$. Strict inequality and unrestricted attainment persist for all
corner data in a sufficiently small open neighbourhood of this example.
\end{theorem}
The proof in \cref{fd:counterexample-proof} supplies a nonnegative
barycentric identity, scalar certificates excluding every pencil
parameter for both families, and a compactness argument excluding
degenerate limits. Its transversality value is $3/2$.
A second rational instance has a larger certified orbit gap; its exact
construction and qualification to $\famA$ are in \cref{fd:second-example}.

The tangent-edge mechanism does not characterize all failures.
\begin{proposition}[Support one and transverse edges do not suffice]
\label{fd:support-one}
For the same $S$ and costs, take $\sbar=(-2,3,2)$ and rays to
\[
 v_1=(0,0,0),\qquad v_2=(6,-2,1/4),\qquad v_3=(1,-5/2,1/2).
\]
Then $\zK=1$, attained only at $v_1$. All three simplex edges at $v_1$
are transverse, with positive derivative values $2,1/4,1/2$.
Nevertheless $\zA<1$: no nonzero $F$ satisfies the closed PSD conditions
at all four simplex vertices. The inactive-ray KKT multipliers are
$1/8$ and $1/4$.
\end{proposition}
\Cref{fd:support-one-proof} gives a nonnegative barycentric identity
and a small integer dual certificate. The full simplex geometry matters
even when the minimizing ray and the inactive-ray multipliers are regular.

\begin{proposition}[A maximal polyhedral set outside the orbit completion]
\label{fd:wedge}
For any $a,b\in\R$, the set
\[
 W=\{(x,y,w):x\ge a,\ y\le b,\ w\ge bx+ay-ab\}
\]
is maximal $\{w\le xy\}$-free and belongs to neither transformation
family. More generally, neither family contains two distinct feasible
boundary contacts on a common bilinear ruling.
\end{proposition}
\Cref{fd:wedge-proof} proves maximality by explicit interior mixtures
and exclusion by the rank-one contact identity. This illustrates a
missing maximal set directly, while \cref{fd:counterexample,fd:support-one}
show that missing sets can affect an objective bound.
